% !TEX program = xelatex
\documentclass[12pt,a4paper]{article}

\usepackage{fontspec}
\usepackage{polyglossia}
\setmainlanguage{russian}
\setotherlanguage{english}
\setmainfont{DejaVu Serif}
\setsansfont{DejaVu Sans}
\setmonofont{DejaVu Sans Mono}

\usepackage{amsmath,amssymb,amsfonts,mathtools,bm}
\usepackage{geometry}
\usepackage{booktabs}
\usepackage{longtable}
\usepackage{array}
\usepackage{enumitem}
\usepackage{hyperref}
\usepackage{microtype}
\geometry{margin=2.5cm}
\hypersetup{colorlinks=true,linkcolor=black,urlcolor=blue,citecolor=black}
\setlength{\emergencystretch}{3em}
\sloppy

\DeclareMathOperator{\Adm}{Adm}
\DeclareMathOperator{\Struct}{Struct}
\DeclareMathOperator{\Dist}{Dist}
\DeclareMathOperator{\Dom}{Dom}
\DeclareMathOperator{\Cond}{Cond}
\DeclareMathOperator{\Loss}{Loss}
\DeclareMathOperator{\Tr}{Tr}
\DeclareMathOperator{\Diff}{Diff}
\DeclareMathOperator{\Spin}{Spin}
\DeclareMathOperator{\Var}{Var}
\DeclareMathOperator{\Spec}{Spec}
\DeclareMathOperator{\Lie}{Lie}

\newcommand{\M}{\mathcal M_U}
\newcommand{\MI}{\mathcal M_I}
\newcommand{\EI}{\mathcal E_I}
\newcommand{\Hilb}{\mathcal H}
\newcommand{\Prob}{\mathbb P}
\newcommand{\Lag}{\mathcal L}
\newcommand{\Acal}{\mathbb A_I}
\newcommand{\Fcal}{\mathbb F_I}
\newcommand{\TMI}{\mathcal T_{\mathrm{MI}}}
\newcommand{\eps}{\varepsilon}
\newcommand{\dd}{\,d}
\newcommand{\Rplus}{\mathbb R_+}
\newcommand{\gfrak}{\mathfrak G}
\newcommand{\statementbox}[1]{\begin{center}\fbox{\begin{minipage}{0.9\linewidth}\centering #1\end{minipage}}\end{center}}

\title{\textbf{Теория многообразных интерфейсов}\\
\large Регуляризованный эффективный физический слой и интерфейсно-полевая программа}
\author{Salkutsan Aleksey}
\date{Май 2026}

\begin{document}
\maketitle

\begin{abstract}
В работе представлена публикационная версия физического слоя Теории многообразных интерфейсов как регуляризованной эффективной интерфейсно-полевой исследовательской программы. Центральная формула состоит в том, что геометрия задаёт область допустимых физических переходов, квантовая теория задаёт распределение по этой области, а интерфейсное поле осуществляет переход от распределённых возможностей к устойчивым физическим структурам. Формально это записывается как
\[
 g\Rightarrow\Adm_c(g)\Rightarrow\Hilb_I(g)\Rightarrow\Psi_I
 \Rightarrow\Dist_{\mathrm{poss}}(\Psi_I;g)\Rightarrow\chi_I\Rightarrow\Struct_I.
\]
Вводятся минимальное интерфейсное поле $\chi_I$, эффективное действие, регуляризованное пространство геометрий, интерфейсное расстояние между классами геометрий, физическая нормировка коэффициента интерфейсной декогеренции, условия восстановления общей теории относительности и Стандартной модели, а также наблюдаемая матрица для квантово-геометрической декогеренции, тёмного сектора, чёрных дыр и интерфейсно-полевой архитектуры объединения взаимодействий. Работа не утверждает, что уже построена завершённая фундаментальная теория или окончательная Theory of Everything. Её результатом является строгая эффективная модель с явными определениями, пределами применимости, проверяемыми гипотезами и открытыми задачами.
\end{abstract}

\noindent\textbf{Ключевые слова:} Теория многообразных интерфейсов; интерфейсное поле; допустимые переходы; квантовая геометрия; декогеренция; тёмный сектор; чёрные дыры; Стандартная модель; объединение взаимодействий; эффективная теория поля.

\tableofcontents
\newpage

\section{Введение}
Фундаментальная физика содержит несколько мощных, но концептуально раздельных описаний. Общая теория относительности описывает пространство-время через метрику $g_{\mu\nu}$ и связывает её с энергией и импульсом через уравнение Эйнштейна. Квантовая теория описывает физическое состояние как вектор, функционал или матрицу плотности, задающие распределение возможных исходов. Стандартная модель описывает калибровочные взаимодействия и фермионные поля. Космология и физика чёрных дыр добавляют собственные проблемы: тёмный сектор, горизонт, энтропию, информационный парадокс и квантовую природу геометрии.

Теория многообразных интерфейсов предлагает общий язык для связи трёх операций: допустимости, распределения и структурирования. В физическом слое ТМИ геометрия отвечает на вопрос, какие переходы физически возможны; квантовая теория отвечает на вопрос, как распределены возможности; интерфейсное поле отвечает на вопрос, как распределение возможностей переходит в устойчивую структуру.

Главная цепочка работы:
\[
\boxed{
 g\Rightarrow\Adm_c(g)\Rightarrow\Hilb_I(g)\Rightarrow\Psi_I
 \Rightarrow\Dist_{\mathrm{poss}}(\Psi_I;g)\Rightarrow\chi_I\Rightarrow\Struct_I .}
\]
Смысл этой цепочки:
\[
\boxed{\text{геометрия задаёт допустимость;}\qquad
\text{квантовая теория задаёт распределение;}\qquad
\text{интерфейсное поле задаёт структурирование.}}
\]

Цель статьи — собрать эту цепочку в цельный физический слой: дать определения, записать минимальное действие, вывести уравнения движения, показать основные режимы, определить наблюдаемые величины, зафиксировать ограничения и сформулировать статус модели. Важная методологическая позиция состоит в том, что ТМИ-Physics рассматривается как эффективная интерфейсно-полевая программа, а не как завершённая фундаментальная теория. Это защищает текст от чрезмерных утверждений и делает его пригодным для дальнейшего усиления.

\section{Связь с предыдущими материалами ТМИ}
Настоящая публикация является синтетической ступенью физической ветки Теории многообразных интерфейсов. Предыдущие модули рассматривали математическую абстракцию ТМИ, осторожный интерфейсный подход к квантовой гравитации, квантово-геометрическую декогеренцию, тёмный сектор как интерфейсное поле, чёрные дыры как причинно-информационные интерфейсы и интерфейсно-полевую архитектуру объединения фундаментальных взаимодействий. В данной работе эти направления не трактуются как разрозненные ветки. Они объединяются как режимы одной минимальной системы:
\[
\boxed{(g_{\mu\nu},\Psi,\chi_I,\Acal,\mathcal G_I)\Rightarrow(QG,Dark,BH,SM,Gen,ToE).}
\]
Здесь $QG$ обозначает квантово-геометрический режим, $Dark$ — тёмный сектор, $BH$ — чёрные дыры, $SM$ — Стандартную модель как предельный сектор, $Gen$ — поколенческую структуру фермионов, а $ToE$ — программу объединения через универсальную интерфейсную связность.

\section{Геометрическая допустимость физических переходов}
Пусть $(\M,g)$ — пространство-время с лоренцевой метрикой. В каждой точке $x\in\M$ задано касательное пространство $T_x\M$. Будущий интерфейсный причинный конус определяется как
\[
\boxed{
\mathcal C^{+}_{I,g}(x):=
\{v\in T_x\M\setminus\{0\}\mid g_x(v,v)\leq0,\ v\ \text{направлен в будущее}\}.
}
\]
Это не новый объект сверх общей теории относительности. Это стандартный будущий причинный конус, записанный как множество локально допустимых направлений физического интерфейсного перехода:
\[
\boxed{\mathcal C^{+}_{I,g}(x)\equiv\mathcal C^{+}_{g}(x).}
\]
Его граница задаётся условием
\[
 g_x(v,v)=0,
\]
поэтому
\[
\boxed{\partial\mathcal C^{+}_{I,g}(x)\equiv\mathcal C^{+}_{\mathrm{light},g}(x).}
\]
Следовательно, световой конус в ТМИ интерпретируется как граница допустимого причинного интерфейсного перехода.

Пусть $J_g^+(x)$ — причинное будущее точки $x$. Тогда множество причинно допустимых переходов определяется так:
\[
\boxed{
\Adm_c(g):=\{(x,y)\in\M\times\M\mid y\in J_g^+(x)\}.
}
\]
Эта формула фиксирует первый слой теории:
\[
\boxed{g\Rightarrow\mathcal C^+_{I,g}(x)\Rightarrow J_g^+(x)\Rightarrow\Adm_c(g).}
\]
Физический смысл: метрика задаёт не только расстояния и кривизну, но и структуру физически допустимых переходов.

\section{Квантовое распределение по области допустимых переходов}
Геометрия задаёт область допустимого, но сама по себе не говорит, какой из допустимых переходов будет реализован. Поэтому вводится квантовый слой. Рассмотрим измеримое пространство
\[
(\Adm_c(g),\mathcal B_g,\mu_g),
\]
где $\mathcal B_g$ — $\sigma$-алгебра измеримых подмножеств, а $\mu_g$ — эффективная или регуляризованная мера. Полная фундаментальная мера на общем пространстве допустимых переходов или геометрий в данной работе не считается построенной. Допускаются четыре уровня:
\[
\boxed{\mu_g\in\{\mu_g^{disc},\mu_g^{mini},\mu_g^{cyl},\mu_g^{PI}\}.}
\]
Здесь $\mu_g^{disc}$ — дискретная мера, $\mu_g^{mini}$ — минисуперпространственная мера, $\mu_g^{cyl}$ — цилиндрическая мера на регуляризованной системе приближений, $\mu_g^{PI}$ — эффективная мера интеграла по путям.

Гильбертово пространство интерфейсных возможностей:
\[
\boxed{
\Hilb_I(g):=L^2(\Adm_c(g),\mathcal B_g,\mu_g).
}
\]
Если $\psi_I\in\Hilb_I(g)$, то вероятность на множестве допустимых переходов задаётся мерой
\[
\boxed{
 d\Prob_{\Psi,g}(\Phi)=|\psi_I(\Phi)|^2\,d\mu_g(\Phi),\qquad \Phi\in\Adm_c(g).
}
\]
Тем самым второй слой записывается как
\[
\boxed{\Adm_c(g)\Rightarrow\Hilb_I(g)\Rightarrow\Psi_I\Rightarrow\Dist_{\mathrm{poss}}(\Psi_I;g).}
\]

\section{Интерфейсное структурирование возможностей}
Распределение возможностей не тождественно реализованной структуре. Для перехода от распределения к устойчивому физическому режиму вводится интерфейсное поле. В минимальной модели
\[
\boxed{\chi_I:\M\to\mathbb R,}
\]
а в более общей геометрической форме
\[
\boxed{\chi_I\in\Gamma(\mathcal E_\chi),\qquad \mathcal E_\chi\to\M.}
\]
Интерфейсное структурирование определяется как отображение
\[
\boxed{
\mathcal S_I:\Dist_{\mathrm{poss}}(\Psi_I;g)\longrightarrow\Struct_I.
}
\]
Более полно:
\[
\boxed{
\Struct_I=\mathcal S_I(g,\Psi_I,\chi_I,\Adm_c(g)).
}
\]
В простейшей модели двух геометрических альтернатив
\[
 |\Psi_g\rangle=c_1|g_1\rangle+c_2|g_2\rangle,
 \qquad |c_1|^2+|c_2|^2=1,
\]
когерентность между альтернативами задаётся недиагональным элементом матрицы плотности. Интерфейсное структурирование описывается законом
\[
\boxed{
 |\rho_{12}(t)|=e^{-\Gamma_I t}|\rho_{12}(0)|.
}
\]
Тогда естественная мера структурирования:
\[
\boxed{
\Struct_I(t):=1-\frac{|\rho_{12}(t)|}{|\rho_{12}(0)|}=1-e^{-\Gamma_I t}.
}
\]

\section{Минимальное эффективное действие}
Физический слой ТМИ задаётся минимальным эффективным действием
\[
\boxed{S_{TMI}^{min}=S_{GR}+S_{SM}+S_\chi+S_{mix}+S_{dec}^{eff}.}
\]
В локальной форме:
\[
\boxed{
S_{TMI}^{min}=\int_{\M}\sqrt{-g}\left[
\frac{c^4}{16\pi G}(R-2\Lambda)+\Lag_{SM}+\Lag_\chi+\Lag_{mix}
\right]d^4x+S_{dec}^{eff}.
}
\]
Здесь $S_{dec}^{eff}$ не является фундаментальным локальным действием в обычном смысле. Это эффективный член открытой квантовой динамики, отвечающий за интерфейсную декогеренцию.

Минимальный лагранжиан интерфейсного поля:
\[
\boxed{
\Lag_\chi=-\frac12 g^{\mu\nu}\nabla_\mu\chi_I\nabla_\nu\chi_I
-V(\chi_I)-\frac12\xi_I R\chi_I^2.
}
\]
Потенциал:
\[
\boxed{
V(\chi_I)=V_0+\frac12m_I^2\chi_I^2+\frac{\beta_I}{4}\chi_I^4.
}
\]
Устойчивость минимальной модели требует, чтобы потенциал был ограничен снизу. Достаточное простое условие:
\[
\boxed{\beta_I>0,\qquad V(\chi_I)\geq0\ \text{в физически допустимой области}.}
\]

Допустимые смешивания должны сохранять калибровочную структуру Стандартной модели. Поэтому в минимальной версии $\chi_I$ выбирается калибровочным синглетом:
\[
\boxed{\chi_I\sim(\mathbf1,\mathbf1,0)\quad\text{относительно }SU(3)\times SU(2)\times U(1)_Y.}
\]
Пример калибровочно безопасного смешивания:
\[
\boxed{
\Lag_{mix}^{gauge}=-\frac14\sum_a\zeta_a\frac{\chi_I^2}{M_I^2}F^a_{\mu\nu}F_a^{\mu\nu}.
}
\]

\section{Уравнения движения минимальной модели}
Варьирование по метрике даёт обобщённое уравнение геометрии:
\[
\boxed{
G_{\mu\nu}+\Lambda g_{\mu\nu}
=\frac{8\pi G}{c^4}\left(T_{\mu\nu}^{SM}+T_{\mu\nu}^{\chi}+T_{\mu\nu}^{mix}\right)+\mathcal D_{\mu\nu}^{eff}.
}
\]
Член $\mathcal D_{\mu\nu}^{eff}$ обозначает возможную эффективную поправку от открытой квантовой динамики. Если обратное влияние декогеренции на геометрию не рассматривается, можно положить $\mathcal D_{\mu\nu}^{eff}=0$.

Варьирование по $\chi_I$ даёт уравнение интерфейсного поля:
\[
\boxed{
\Box_g\chi_I-(m_I^2+\xi_I R)\chi_I-\beta_I\chi_I^3=-J_I.
}
\]
Источник $J_I$ кодирует возбуждение интерфейсного поля квантово-геометрическими или материальными конфигурациями. В QG-режиме минимально:
\[
\boxed{J_I=\alpha_I |c_1|^2|c_2|^2d_I^2(g_1,g_2).}
\]
Квантовая открытая динамика записывается как
\[
\boxed{
\frac{d\rho}{dt}=-\frac{i}{\hbar}[\widehat H_{eff},\rho]+\mathcal L_{dec}^{I}[\rho].
}
\]
Минимальный декогеренционный супероператор:
\[
\boxed{
\mathcal L_I[\rho]=-\frac{1}{2}\kappa_I[\widehat D_I,[\widehat D_I,\rho]].
}
\]
Калибровочные поля получают модифицированные уравнения вида
\[
\boxed{
D_\mu\left[\left(1+\zeta_a\frac{\chi_I^2}{M_I^2}\right)F_a^{\mu\nu}\right]=J_a^\nu.
}
\]

\section{Регуляризованное пространство геометрий}
Для квантово-геометрического режима нужно перейти от конечной суперпозиции геометрий к пространству геометрических возможностей. Пусть $\mathcal G(\M)$ — пространство допустимых лоренцевых метрик. Физическая геометрия задаётся не координатной записью метрики, а классом по диффеоморфизмам:
\[
\boxed{[g]=\{\varphi^*g\mid \varphi\in\Diff(\M)\}.}
\]
Пространство физических геометрий:
\[
\boxed{\gfrak:=\mathcal G(\M)/\Diff(\M).}
\]
Полная мера на $\gfrak$ в данной работе не считается построенной. Вместо этого вводится конечномерная регуляризация:
\[
\boxed{\mathcal G_N\subset\mathcal G(\M),\qquad \gfrak_N=\mathcal G_N/\Diff_N.}
\]
Рабочее гильбертово пространство геометрий:
\[
\boxed{
\Hilb_G^{eff}=L^2(\gfrak_N,\mathcal B_N,\mu_G^{cyl}).
}
\]
Здесь $\mu_G^{cyl}$ — регуляризованная цилиндрическая мера. Она не объявляется окончательной фундаментальной мерой, но даёт вычислимую иерархию приближений.

Интерфейсное расстояние между физическими геометриями должно сравнивать классы, а не координатные записи:
\[
\boxed{d_I([g_1],[g_2]).}
\]
Общая эффективная форма:
\[
\boxed{
d_I([g_1],[g_2])=\frac{1}{\mathcal N_I}\inf_{\varphi\in\Diff(\M)}
\left[\int_\Omega w_I(x)\|g_1-\varphi^*g_2\|_{g_0}^2dV_{g_0}\right]^{1/2}.
}
\]
В минисуперпространственной модели с масштабным фактором $a$:
\[
\boxed{d_I(g(a_1),g(a_2))=\left|\ln\frac{a_1}{a_2}\right|.}
\]

\section{Физическая нормировка коэффициента декогеренции}
QG-модуль использует наблюдаемую добавку
\[
\boxed{\Delta\Gamma_I=\kappa_I x_I,}
\]
где
\[
\boxed{x_I=|c_1|^2|c_2|^2d_{I,N}^2([g_1],[g_2]).}
\]
Если $x_I$ безразмерен, то $\kappa_I$ имеет размерность обратного времени. Поэтому физически корректная нормировка задаётся через характерную частоту интерфейсного поля:
\[
\boxed{\omega_I=\frac{m_Ic^2}{\hbar}.}
\]
Тогда
\[
\boxed{\kappa_I=\frac{m_Ic^2}{\hbar}\mathcal K_I.}
\]
В минимальной эффективной модели
\[
\boxed{\mathcal K_I=\lambda_I\alpha_I^2,}
\]
а более общая запись допускает поправки от кривизны, самодействия и энергетического масштаба:
\[
\boxed{
\kappa_I=\frac{m_Ic^2}{\hbar}\lambda_I\alpha_I^2\,
\mathcal F_I\left(\frac{\xi_I R}{m_I^2},\beta_I,\frac{E}{M_I},\ldots\right).
}
\]
Таким образом, $\kappa_I$ уже не является полностью произвольным параметром: он связан с масштабом интерфейсного поля и его безразмерными связями.

\section{Полная квантовая динамика метрики в регуляризованной модели}
Состояние квантовой геометрии теперь рассматривается как функция на пространстве физических геометрий:
\[
\boxed{\Psi_G([g],t)\in\Hilb_G^{eff}.}
\]
Нормировка:
\[
\boxed{\int_{\gfrak_N}|\Psi_G([g],t)|^2d\mu_G^{cyl}([g])=1.}
\]
Эффективная динамика может быть записана как
\[
\boxed{i\hbar\frac{\partial}{\partial t}\Psi_G([g],t)=\widehat H_G^{eff}\Psi_G([g],t).}
\]
Однако параметр $t$ должен пониматься осторожно: это эффективное или внутреннее время выбранной регуляризации, а не полное решение проблемы времени в квантовой гравитации. Более ковариантная форма задаётся ограничением
\[
\boxed{\widehat{\mathcal H}_{TMI}\Psi=0.}
\]
Для открытой динамики геометрии удобнее использовать матрицу плотности:
\[
\boxed{
\frac{d\rho_G}{dt}=-\frac{i}{\hbar}[\widehat H_G^{eff},\rho_G]
-\frac12\kappa_I[\widehat D_I,[\widehat D_I,\rho_G]].
}
\]
Если $\widehat D_I$ различает геометрии через $d_{I,N}$, то
\[
\boxed{\Gamma_I([g_1],[g_2])=\kappa_I d_{I,N}^2([g_1],[g_2]).}
\]
Следовательно,
\[
\boxed{\rho_G([g_1],[g_2];t)=e^{-\Gamma_I([g_1],[g_2])t}\rho_G([g_1],[g_2];0).}
\]
Классическая метрика возникает как декогерированный предел квантового распределения геометрий:
\[
\boxed{\rho_G([g_1],[g_2])\to0\quad([g_1]\ne[g_2])\Rightarrow g_{\mu\nu}^{cl}.}
\]

\section{Восстановление Стандартной модели и проекции взаимодействий}
Безопасная структура универсальной интерфейсной группы задаётся как
\[
\boxed{G_I=G_{\mathrm{grav}}\ltimes G_{SM}\ltimes G_\chi,}
\]
где
\[
\boxed{G_{SM}=SU(3)\times SU(2)\times U(1)_Y,\qquad G_{\mathrm{grav}}\supset\Spin(1,3).}
\]
Алгебра минимальной модели:
\[
\boxed{\mathfrak g_I=\mathfrak g_{\mathrm{grav}}\oplus\mathfrak{su}(3)\oplus\mathfrak{su}(2)\oplus\mathfrak u(1)\oplus\mathfrak g_\chi.}
\]
Универсальная интерфейсная связность:
\[
\boxed{\Acal=\omega^{ab}J_{ab}+G^AT_A+W^iT_i+BY+A_\chi^r\Xi_r.}
\]
Универсальная кривизна:
\[
\boxed{\Fcal=d\Acal+\Acal\wedge\Acal.}
\]
Вводятся проекции алгебры
\[
\boxed{\pi_a:\mathfrak g_I\to\mathfrak g_a,
\qquad a\in\{G,3,2,1,\chi\}.}
\]
Они удовлетворяют
\[
\boxed{\pi_a^2=\pi_a,\\ \pi_a\pi_b=0\ (a\ne b),\qquad \sum_a\pi_a=\mathrm{id}_{\mathfrak g_I}.}
\]
Тогда
\[
\boxed{F_a=\pi_a(\Fcal).}
\]
Теперь утверждение, что взаимодействия являются проекциями универсальной интерфейсной кривизны, имеет строгий смысл.

Электромагнитный сектор возникает после электрослабого нарушения симметрии:
\[
\boxed{SU(2)\times U(1)_Y\to U(1)_{EM}.}
\]
Поэтому
\[
\boxed{A_\mu^{EM}=\sin\theta_W W_\mu^3+\cos\theta_WB_\mu.}
\]
Безопасный предел Стандартной модели:
\[
\boxed{g_{\mu\nu}=g_{\mu\nu}^{fixed},\quad \chi_I\to0,
\quad S_{mix}\to0\Rightarrow S_{TMI}^{min}\to S_{SM}.}
\]
С гравитацией:
\[
\boxed{\chi_I\to0,
\quad S_{mix}\to0\Rightarrow S_{TMI}^{min}\to S_{GR}+S_{SM}.}
\]

\section{Фермионные поколения как устойчивые интерфейсные моды}
Стандартная модель содержит три поколения фермионов, но не выводит их число из первых принципов. В ТМИ вводится внутреннее поколенческое пространство $\Hilb_{gen}$ и оператор интерфейсного структурирования поколений
\[
\boxed{\mathcal G_I:\Hilb_{gen}\to\Hilb_{gen}.}
\]
Спектральная задача:
\[
\boxed{\mathcal G_I\psi_f=\lambda_f\psi_f.}
\]
Фермионное состояние имеет вид
\[
\boxed{\Psi_f\in\Hilb_{SM}\otimes\Hilb_{gen},
\qquad \Psi_f=\psi_{SM}\otimes\psi_{gen}^{(f)}.}
\]
Поколенческий оператор не должен разрушать калибровочную структуру:
\[
\boxed{[\mathcal G_I,G_{SM}]=0.}
\]
Физические поколения отождествляются с устойчивыми модами:
\[
\boxed{f\leftrightarrow\psi_f\in\Spec_{stable}(\mathcal G_I).}
\]
Слабая форма результата: поколения формализованы как устойчивые интерфейсные моды. Сильная форма остаётся открытой:
\[
\boxed{|\Spec_{stable}(\mathcal G_I)|=3.}
\]
Массы могут зависеть от интерфейсной моды:
\[
\boxed{m_f=\frac{v_H}{\sqrt2}\left[y_f^{SM}+\eta_f\frac{\lambda_f\langle\chi_I\rangle}{M_I}\right].}
\]
Смешивание поколений может возникать из недиагональности поколенческого оператора:
\[
\boxed{\Delta Y_{ij}^{I}=\eta_{ij}\frac{\langle\psi_i|\mathcal G_I|\psi_j\rangle}{M_I}\langle\chi_I\rangle.}
\]

\section{Физические режимы минимальной системы}
Минимальная система содержит четыре ключевых объекта:
\[
\boxed{g_{\mu\nu},\qquad\Psi,
\qquad\chi_I,
\qquad\Acal.}
\]
Из них возникают основные режимы.

\subsection{Квантово-геометрический режим}
QG-режим описывает подавление когерентности между различными геометрическими альтернативами:
\[
\boxed{\Delta\Gamma_I=\kappa_Ix_I,
\qquad x_I=|c_1|^2|c_2|^2d_I^2([g_1],[g_2]).}
\]
Наблюдаемый эффект — избыточная геометрическая декогеренция.

\subsection{Тёмный сектор}
Тёмный сектор описывается как эффективный вклад интерфейсного поля:
\[
\boxed{T_{\mu\nu}^{dark}=T_{\mu\nu}^{I}[\chi_I,g].}
\]
Разложение
\[
\boxed{\chi_I(x,t)=\chi_0(t)+\delta\chi(x,t)}
\]
задаёт два режима:
\[
\boxed{\chi_0(t)\Rightarrow DE,
\qquad \delta\chi(x,t)\Rightarrow DM.}
\]
Для фонового режима
\[
\boxed{w_I=\frac{\frac12\dot\chi_0^2-V(\chi_0)}{\frac12\dot\chi_0^2+V(\chi_0)}.}
\]
Если $V(\chi_0)\gg\frac12\dot\chi_0^2$, то $w_I\approx-1$.

\subsection{Чёрные дыры}
В ТМИ чёрная дыра — это односторонний причинно-информационный интерфейс. Для чёрной дыры Шварцшильда
\[
\boxed{r_s=\frac{2GM}{c^2}.}
\]
Горизонт $r=r_s$ является границей множества допустимых внешних причинных переходов:
\[
\boxed{\mathcal H_g=\partial\Adm_{out}.}
\]
Внутри горизонта классический выходной канал запрещён:
\[
\boxed{r<r_s\Rightarrow\Adm_{out}^{class}=\varnothing.}
\]
Квантовый внешний канал связан с излучением Хокинга. Температура
\[
\boxed{T_H=\frac{\hbar c^3}{8\pi GMk_B}}
\]
интерпретируется как интенсивность квантового внешнего интерфейсного канала.

\subsection{Архитектура объединения}
ToE-режим в данной работе понимается осторожно: это не завершённая теория всего, а интерфейсно-полевая программа объединения. Четыре взаимодействия описываются как проекции универсальной интерфейсной кривизны:
\[
\boxed{\Fcal\Rightarrow(F_G,F_3,F_2,F_1,F_\chi),
\qquad F_a=\pi_a(\Fcal).}
\]
Интерфейсные поправки к калибровочным связям имеют вид
\[
\boxed{\frac{\Delta g_a}{g_a}\approx-\frac12\zeta_a\frac{\chi_I^2}{M_I^2}.}
\]

\section{Наблюдаемая матрица и критерии фальсификации}
Физический слой получает смысл проверяемой программы только при наличии наблюдаемых выходов. Вводится матрица
\[
\boxed{
\mathcal T_{TMI}^{phys}=\left(\Delta\Gamma_I,
 w_I(z),G_{eff}(k,z),\rho_{corr},\Delta_a(E)\right).
}
\]
Здесь:
\begin{itemize}[leftmargin=2em]
\item $\Delta\Gamma_I=\Gamma_{obs}-\Gamma_{std}$ — избыточная геометрическая декогеренция;
\item $w_I(z)$ — возможное отклонение тёмной энергии от идеального $w=-1$;
\item $G_{eff}(k,z)$ — возможная эффективная гравитационная связь в тёмном секторе;
\item $\rho_{corr}$ — корреляционный вклад в излучение чёрной дыры;
\item $\Delta_a(E)$ — интерфейсная поправка к калибровочным связям.
\end{itemize}

Для QG-режима:
\[
\boxed{H_0^{QG}:\Delta\Gamma_I=0,}
\qquad
\boxed{H_{TMI}^{QG}:\Delta\Gamma_I=\kappa_Ix_I.}
\]
Для тёмного сектора:
\[
\boxed{H_0^{Dark}:w=-1,\ DM\ \text{и}\ DE\ \text{независимы},}
\]
\[
\boxed{H_{TMI}^{Dark}:DM\ \text{и}\ DE\ \text{являются режимами одного поля}\ \chi_I.}
\]
Для чёрных дыр:
\[
\boxed{H_0^{BH}:\rho_{rad}=\rho_{thermal},}
\qquad
\boxed{H_{TMI}^{BH}:\rho_{rad}=\rho_{thermal}+\rho_{corr}.}
\]
Для объединения:
\[
\boxed{H_0^{ToE}:\Delta_a(E)=0,}
\qquad
\boxed{H_{TMI}^{ToE}:\Delta_a(E)=\zeta_a\frac{\chi_I^2(E)}{M_I^2}.}
\]

\section{Условия согласованности}
Физический слой допустим только при восстановлении известных теорий в соответствующих пределах.

\subsection{Предел общей теории относительности}
Если
\[
\boxed{\chi_I\to0,
\quad S_{mix}\to0,
\quad S_{dec}^{eff}\to0,}
\]
то
\[
\boxed{S_{TMI}^{min}\to S_{GR}+S_{matter}.}
\]
При отсутствии материи:
\[
\boxed{S_{TMI}^{min}\to S_{GR}.}
\]

\subsection{Квантовый предел}
При фиксированной геометрии и выключенной интерфейсной декогеренции:
\[
\boxed{g_{\mu\nu}=g_{\mu\nu}^{fixed},\quad\Gamma_I\to0,}
\]
должна восстанавливаться стандартная квантовая динамика:
\[
\boxed{i\hbar\frac{\partial}{\partial t}\Psi=\widehat H\Psi.}
\]

\subsection{Предел Стандартной модели}
Если
\[
\boxed{g_{\mu\nu}=g_{\mu\nu}^{fixed},\quad \chi_I\to0,
\quad \zeta_a,\eta_H,y_{I,f}\to0,}
\]
то
\[
\boxed{S_{TMI}^{min}\to S_{SM}.}
\]

\section{Уровни утверждений и ограничения}
Чтобы не смешивать разные типы утверждений, вводятся три уровня.

\subsection{Уровень A: внутренние следствия}
Это утверждения, следующие из определений. Например,
\[
\boxed{g\Rightarrow\Adm_c(g)}
\]
следует после определения $\Adm_c(g)$ через $J_g^+(x)$. Аналогично,
\[
\boxed{\psi_I\in L^2(\Adm_c(g),\mathcal B_g,\mu_g)
\Rightarrow d\Prob_{\Psi,g}=|\psi_I|^2d\mu_g.}
\]

\subsection{Уровень B: феноменологические гипотезы}
Это проверяемые физические предположения:
\[
\boxed{\Gamma_I=\kappa_I|c_1|^2|c_2|^2d_I^2([g_1],[g_2]),}
\]
\[
\boxed{T_{\mu\nu}^{dark}\sim T_{\mu\nu}^{\chi},}
\qquad
\boxed{\Delta_a(E)=\zeta_a\frac{\chi_I^2(E)}{M_I^2}.}
\]
Их нельзя называть доказанными фактами природы.

\subsection{Уровень C: открытые задачи}
Открытыми остаются: полная фундаментальная мера $\mu_G$ в континуальном пределе, строгий переход $N\to\infty$, фундаментальный оператор $\widehat{\mathcal H}_{TMI}$, полное решение проблемы времени, вывод $SU(3)\times SU(2)\times U(1)$ как единственного предела, вывод ровно трёх поколений, численные массы, CKM и PMNS, а также связь с конкретными экспериментальными ограничениями.

\section{Итоговая теорема регуляризованного физического слоя}
\textbf{Теорема.} Пусть задано регуляризованное пространство физических геометрий
\[
\boxed{\gfrak_N=\mathcal G_N/\Diff_N}
\]
с мерой $\mu_G^{cyl}$ и интерфейсным расстоянием $d_{I,N}$. Пусть
\[
\boxed{\Hilb_G^{eff}=L^2(\gfrak_N,\mathcal B_N,\mu_G^{cyl})}
\]
и коэффициент интерфейсной декогеренции имеет форму
\[
\boxed{\kappa_I=\frac{m_Ic^2}{\hbar}\mathcal K_I.}
\]
Если открытая динамика геометрии задаётся уравнением
\[
\boxed{
\frac{d\rho_G}{dt}=-\frac{i}{\hbar}[\widehat H_G^{eff},\rho_G]
-\frac12\kappa_I[\widehat D_I,[\widehat D_I,\rho_G]],
}
\]
то недиагональные элементы между различными геометрическими классами подавляются со скоростью
\[
\boxed{\Gamma_I([g_1],[g_2])=\kappa_Id_{I,N}^2([g_1],[g_2]).}
\]
Следовательно,
\[
\boxed{\rho_G([g_1],[g_2])\to0\quad([g_1]\ne[g_2])}
\]
и возникает эффективная классическая геометрия $g_{\mu\nu}^{cl}$.

\textbf{Смысл теоремы:}
\[
\boxed{\text{классическая геометрия есть интерфейсно-декогерированный предел квантового распределения геометрий.}}
\]

\section{Заключение}
Физический слой Теории многообразных интерфейсов в данной работе доведён до статуса регуляризованной эффективной интерфейсно-полевой модели. Его центральная формула:
\[
\boxed{
 g\Rightarrow\Adm_c(g)\Rightarrow\Hilb_I(g)\Rightarrow\Psi_I
 \Rightarrow\Dist_{\mathrm{poss}}(\Psi_I;g)\Rightarrow\chi_I\Rightarrow\Struct_I.
}
\]
Расширенная архитектура:
\[
\boxed{(g_{\mu\nu},\Psi,\chi_I,\Acal,\mathcal G_I)
\Rightarrow(QG,SM,Gen,Dark,BH,ToE).}
\]
В этой форме ТМИ-Physics имеет математический каркас, физическое действие, уравнения движения, режимы, предельные случаи и наблюдаемую матрицу. При этом она не является завершённой фундаментальной теорией. Корректный статус:
\[
\boxed{TMI\text{-}Physics=\text{регуляризованная эффективная интерфейсно-полевая исследовательская программа}.}
\]

\appendix
\section{Единый реестр обозначений}
\begin{longtable}{p{0.26\linewidth}p{0.66\linewidth}}
\toprule
Обозначение & Значение \\
\midrule
$\M$ & пространство-время Вселенной \\
$g_{\mu\nu}$ & лоренцева метрика \\
$T_x\M$ & касательное пространство в точке $x$ \\
$\mathcal C^+_{I,g}(x)$ & будущий интерфейсный причинный конус \\
$J_g^+(x)$ & причинное будущее точки $x$ \\
$\Adm_c(g)$ & множество причинно допустимых переходов \\
$\mathcal B_g$ & $\sigma$-алгебра измеримых подмножеств \\
$\mu_g$ & эффективная или регуляризованная мера \\
$\Hilb_I(g)$ & гильбертово пространство интерфейсных возможностей \\
$\Psi_I,\psi_I$ & квантовое состояние и волновая функция интерфейсного слоя \\
$\Dist_{\mathrm{poss}}$ & распределение возможностей \\
$\chi_I$ & интерфейсное поле \\
$\Struct_I$ & структурированный физический режим \\
$\Gamma_I$ & скорость интерфейсной декогеренции \\
$\kappa_I$ & коэффициент силы интерфейсного канала \\
$d_I$ & интерфейсное расстояние между геометриями \\
$\Acal$ & универсальная интерфейсная связность \\
$\Fcal$ & универсальная интерфейсная кривизна \\
$\mathcal G_I$ & поколенческий интерфейсный оператор \\
$\Delta\Gamma_I$ & избыточная интерфейсная декогеренция \\
$w_I(z)$ & параметр состояния интерфейсной тёмной энергии \\
$\rho_{corr}$ & корреляционный вклад в излучение чёрной дыры \\
$\Delta_a(E)$ & интерфейсная поправка к калибровочным связям \\
\bottomrule
\end{longtable}

\section{Логический порядок вывода}
Определения задают $(\M,g)$, $\Adm_c(g)$, $\Hilb_I(g)$, $\Prob_{\Psi,g}$, $\chi_I$ и $\Struct_I$. Постулаты физического слоя:
\[
\boxed{P_1:\ g\Rightarrow\Adm_c(g),}
\qquad
\boxed{P_2:\ \Psi_I\Rightarrow\Dist_{\mathrm{poss}}(\Psi_I;g),}
\qquad
\boxed{P_3:\ \chi_I\Rightarrow\Struct_I.}
\]
Внутренние следствия выводятся из этих определений. Феноменологические гипотезы задают конкретные проверяемые связи: $\Delta\Gamma_I=\kappa_Ix_I$, $T_{\mu\nu}^{dark}\sim T_{\mu\nu}^{\chi}$, $\rho_{rad}=\rho_{thermal}+\rho_{corr}$ и $\Delta_a(E)=\zeta_a\chi_I^2(E)/M_I^2$.

\begin{thebibliography}{99}
\bibitem{TMI-Full} Salkutsan Aleksey, \textit{Theory of Manifold Interfaces: An Axiomatic Conceptual-Mathematical Framework of Interfaciality, Meaning, and Cognition}, Zenodo. \href{https://doi.org/10.5281/zenodo.19940675}{https://doi.org/10.5281/zenodo.19940675}.
\bibitem{TMI-Math} Salkutsan Aleksey, \textit{Mathematical Abstraction of the Theory of Manifold Interfaces}, Zenodo. \href{https://doi.org/10.5281/zenodo.19978895}{https://doi.org/10.5281/zenodo.19978895}.
\bibitem{TMI-CautiousQG} Salkutsan Aleksey, \textit{Theory of Manifold Interfaces. A Cautious Interface Approach to Quantum Gravity}, Zenodo. \href{https://doi.org/10.5281/zenodo.19988505}{https://doi.org/10.5281/zenodo.19988505}.
\bibitem{TMI-QG20} Salkutsan Aleksey, \textit{Theory of Manifold Interfaces TMI-QG-2.0: A Minimal Predictive Model of Interface-Induced Geometric Decoherence}, Zenodo. \href{https://doi.org/10.5281/zenodo.20013976}{https://doi.org/10.5281/zenodo.20013976}.
\bibitem{TMI-BH} Salkutsan Aleksey, \textit{Black Holes in the Theory of Manifold Interfaces}, Zenodo. \href{https://doi.org/10.5281/zenodo.20025226}{https://doi.org/10.5281/zenodo.20025226}.
\bibitem{TMI-Dark} Salkutsan Aleksey, \textit{The Dark Sector as an Interface Field. A Module of the Theory of Manifold Interfaces}, Zenodo. \href{https://doi.org/10.5281/zenodo.20025682}{https://doi.org/10.5281/zenodo.20025682}.
\bibitem{TMI-ToE} Salkutsan Aleksey, \textit{Interface-Field Architecture for the Unification of Fundamental Interactions}, Zenodo. \href{https://doi.org/10.5281/zenodo.20025731}{https://doi.org/10.5281/zenodo.20025731}.
\bibitem{Einstein1915} A. Einstein, \textit{Die Feldgleichungen der Gravitation}, Sitzungsberichte der Preussischen Akademie der Wissenschaften, 1915.
\bibitem{Wald1984} R. M. Wald, \textit{General Relativity}, University of Chicago Press, 1984.
\bibitem{Misner1973} C. W. Misner, K. S. Thorne, J. A. Wheeler, \textit{Gravitation}, W. H. Freeman, 1973.
\bibitem{Weinberg1995} S. Weinberg, \textit{The Quantum Theory of Fields}, Vol. 1, Cambridge University Press, 1995.
\bibitem{Peskin1995} M. E. Peskin and D. V. Schroeder, \textit{An Introduction to Quantum Field Theory}, Addison-Wesley, 1995.
\bibitem{Bekenstein1973} J. D. Bekenstein, \textit{Black holes and entropy}, Physical Review D, 7(8), 1973.
\bibitem{Hawking1975} S. W. Hawking, \textit{Particle creation by black holes}, Communications in Mathematical Physics, 43, 1975.
\bibitem{Zurek2003} W. H. Zurek, \textit{Decoherence, einselection, and the quantum origins of the classical}, Reviews of Modern Physics, 75, 2003.
\bibitem{Schlosshauer2007} M. Schlosshauer, \textit{Decoherence and the Quantum-to-Classical Transition}, Springer, 2007.
\bibitem{Planck2020} Planck Collaboration, \textit{Planck 2018 results. VI. Cosmological parameters}, Astronomy \& Astrophysics, 641, A6, 2020.
\end{thebibliography}

\end{document}
